\documentclass[11pt,a4paper]{article}
\usepackage[utf8]{inputenc}
\usepackage{amsmath,amssymb}
\usepackage{fancyhdr}
\usepackage[margin=2.5cm]{geometry}

\pagestyle{fancy}
\fancyhf{}
\fancyhead[R]{Ivan Mišustin}
\renewcommand{\headrulewidth}{0pt}

% rozšířená matice soustavy (4 sloupce | pravá strana)
\newcommand{\aug}[1]{\left(\begin{array}{rrrr|r} #1 \end{array}\right)}

\begin{document}

\begin{center}
    \LARGE\textbf{Lineární algebra -- domácí úkol č. 2}
\end{center}
\vspace{1em}

Najděte všechna řešení následující soustavy rovnic:
\[
  \begin{array}{rrrrcr}
    3x_1 & {}+ 5x_2 & {}+ 2x_3 & {}+ 3x_4 & = & -3 \\
    2x_1 & {}+ 3x_2 & {}+ 0x_3 & {}+ x_4  & = & -2 \\
     x_1 & {}+ x_2  & {}+ 2x_3 & {}+ x_4  & = & 1 \\
    4x_1 & {}+ 4x_2 & {}+ 3x_3 & {}+ x_4  & = & 2
  \end{array}
\]

Převedu soustavu na rozšířenou matici soustavy:
\[
  \aug{
    3 & 5 & 2 & 3 & -3 \\
    2 & 3 & 0 & 1 & -2 \\
    1 & 1 & 2 & 1 & 1 \\
    4 & 4 & 3 & 1 & 2
  }
\]

A dále pomocí Gaussovy-Jordanovy eliminace převedu tuto matici do redukovaného odstup\-ňova\-ného tvaru:
\begin{align*}
  & \aug{
    3 & 5 & 2 & 3 & -3 \\
    2 & 3 & 0 & 1 & -2 \\
    1 & 1 & 2 & 1 & 1 \\
    4 & 4 & 3 & 1 & 2
  } \sim
  \aug{
    1 & 1 & 2 & 1 & 1 \\
    2 & 3 & 0 & 1 & -2 \\
    3 & 5 & 2 & 3 & -3 \\
    4 & 4 & 3 & 1 & 2
  } \sim
  \aug{
    1 & 1 & 2 & 1 & 1 \\
    0 & 1 & -4 & -1 & -4 \\
    0 & 2 & -4 & 0 & -6 \\
    0 & 0 & -5 & -3 & -2
  } \\[1ex]
  \sim\; & \aug{
    1 & 1 & 2 & 1 & 1 \\
    0 & 1 & -4 & -1 & -4 \\
    0 & 0 & 4 & 2 & 2 \\
    0 & 0 & -5 & -3 & -2
  } \sim
  \aug{
    1 & 1 & 2 & 1 & 1 \\
    0 & 1 & -4 & -1 & -4 \\
    0 & 0 & 2 & 1 & 1 \\
    0 & 0 & -5 & -3 & -2
  } \sim
  \aug{
    1 & 1 & 2 & 1 & 1 \\
    0 & 1 & -4 & -1 & -4 \\
    0 & 0 & 10 & 5 & 5 \\
    0 & 0 & -10 & -6 & -4
  } \\[1ex]
  \sim\; & \aug{
    1 & 1 & 2 & 1 & 1 \\
    0 & 1 & -4 & -1 & -4 \\
    0 & 0 & 10 & 5 & 5 \\
    0 & 0 & 0 & -1 & 1
  } \sim
  \aug{
    1 & 1 & 2 & 1 & 1 \\
    0 & 1 & -4 & -1 & -4 \\
    0 & 0 & 2 & 1 & 1 \\
    0 & 0 & 0 & 1 & -1
  } \sim
  \aug{
    1 & 1 & 2 & 1 & 1 \\
    0 & 1 & -4 & -1 & -4 \\
    0 & 0 & 2 & 0 & 2 \\
    0 & 0 & 0 & 1 & -1
  } \\[1ex]
  \sim\; & \aug{
    1 & 1 & 2 & 1 & 1 \\
    0 & 1 & 0 & -1 & 0 \\
    0 & 0 & 2 & 0 & 2 \\
    0 & 0 & 0 & 1 & -1
  } \sim
  \aug{
    1 & 1 & 2 & 1 & 1 \\
    0 & 1 & 0 & -1 & 0 \\
    0 & 0 & 1 & 0 & 1 \\
    0 & 0 & 0 & 1 & -1
  } \sim
  \aug{
    1 & 1 & 2 & 1 & 1 \\
    0 & 1 & 0 & 0 & -1 \\
    0 & 0 & 1 & 0 & 1 \\
    0 & 0 & 0 & 1 & -1
  } \\[1ex]
  \sim\; & \aug{
    1 & 1 & 2 & 0 & 2 \\
    0 & 1 & 0 & 0 & -1 \\
    0 & 0 & 1 & 0 & 1 \\
    0 & 0 & 0 & 1 & -1
  } \sim
  \aug{
    1 & 1 & 0 & 0 & 0 \\
    0 & 1 & 0 & 0 & -1 \\
    0 & 0 & 1 & 0 & 1 \\
    0 & 0 & 0 & 1 & -1
  } \sim
  \aug{
    1 & 0 & 0 & 0 & 1 \\
    0 & 1 & 0 & 0 & -1 \\
    0 & 0 & 1 & 0 & 1 \\
    0 & 0 & 0 & 1 & -1
  }
\end{align*}

A tedy:
\begin{align*}
  x_1 &= 1 \\
  x_2 &= -1 \\
  x_3 &= 1 \\
  x_4 &= -1
\end{align*}

\end{document}
